
\begin{abstract}
% 150-200 words. A reader should know after the abstract whether the package is for them.
[1 sentence: the problem, in the user's words, not the package's]
[1 sentence: who has this problem (researchers in X, analysts doing Y)]
[1-2 sentences: what \texttt{{{package}}} does, and how it differs from existing tools]
[1 sentence: key design idea (e.g.\ one call, convention over configuration, no dependencies)]
[1 sentence: availability: open source, license, where to get it]

\medskip
\noindent\textbf{Keywords:} R, [domain], [Keyword 3], [Keyword 4], reproducibility
\end{abstract}

\newpage

% =====================================================================
\section{Introduction}\label{sec:intro} % ~10%

\textbf{[The problem]}
[Concrete situation where the problem bites. One example a reader recognises from their own work.]

\textbf{[Why it matters]}
[Cost of the problem: lost time, errors, non-reproducible results. Cite evidence if any exists.]

\textbf{[What this package does]}
[One paragraph: \texttt{{{package}}} provides [main functions], which [what they achieve].]

\textbf{[Roadmap]}
[\Cref{sec:need} states the need, \cref{sec:field} compares existing tools, \cref{sec:design} explains the design, \cref{sec:usage} shows usage, \cref{sec:quality} reports how it is tested, \cref{sec:limits} discusses limitations.]

% =====================================================================
\section{Statement of Need}\label{sec:need} % ~10%
[Who is the target user (field, skill level)? What do they do today without the package? Why is that insufficient? Be specific: a reviewer must be convinced the package fills a real gap, not reinvents a wheel.]

% =====================================================================
\section{State of the Field}\label{sec:field} % ~15%
[Name the existing tools and what each does well \citep{author2020}. Then state precisely what none of them does, which is what this package adds. Be fair to alternatives: say when they are the better choice.]

\begin{table}[htbp]
    \centering
    \small
    \caption{[Comparison with existing tools] Example: droids}
    \label{tab:comparison}
    % example: replace the droids with the tools you compare against
    \begin{tabular}{l*{4}{c}}
        \toprule
        Feature                    & \texttt{{{package}}} & R2-D2      & C-3PO      & BB-8       \\
        \midrule
        Repairs a hyperdrive       & \checkmark           & \checkmark &            &            \\
        Speaks 6 million languages & \checkmark           &            & \checkmark &            \\
        Rolls on sand              & \checkmark           &            &            & \checkmark \\
        Films in \texttt{starwars} & --                   & 7          & 6          & 1          \\
        \bottomrule
    \end{tabular}
\end{table}

% =====================================================================
\section{Design}\label{sec:design} % ~20%

\subsection{Design Principles}
[2-4 principles that explain every decision, e.g.\ \enquote{convention over configuration}, \enquote{every step is idempotent}, \enquote{generated, not hand-written}. One short paragraph each.]

\subsection{Architecture}
[How the functions relate to each other and to what they produce (\cref{fig:architecture}). Which functions are exported, which are internal helpers, and why.]

\begin{figure}[htbp]
    \centering
    \footnotesize
    \begin{tikzpicture}[box/.style={draw, rounded corners, align=center, minimum width=3.2cm, minimum height=1cm},
                        result/.style={draw, dashed, align=center, minimum width=3.2cm, minimum height=1cm}]
        \node[box] (user) at (0,0)  {user call};
        \node[box] (f1)   at (4,1)  {\texttt{[function\_1()]}};
        \node[box] (f2)   at (4,-1) {\texttt{[function\_2()]}};
        \node[result] (o1) at (8.5,1)  {[output 1]};
        \node[result] (o2) at (8.5,-1) {[output 2]};

        \draw[-{Stealth}, thick] (user) -- (f1);
        \draw[-{Stealth}, thick] (f1) -- node[right] {calls} (f2);
        \draw[-{Stealth}, thick] (f1) -- (o1);
        \draw[-{Stealth}, thick] (f2) -- (o2);
    \end{tikzpicture}
    \caption{[Architecture: exported functions and what they produce]}
    \label{fig:architecture}
\end{figure}

\subsection{Key Decisions and Trade-offs}
[What did you choose NOT to do, and why? Each feature you left out is a decision a reviewer will ask about. E.g.\ \enquote{we rely on \texttt{usethis} rather than reimplementing file creation, at the cost of one dependency}.]

% =====================================================================
\section{Usage}\label{sec:usage} % ~20%

\subsection{Installation}

\begin{lstlisting}
remotes::install_github("{{repo}}")
library({{package}})
\end{lstlisting}

\subsection{Minimal Example}
[The smallest call that shows the value of the package. The reader should be able to copy-paste it.]

\begin{lstlisting}
# [what this example does]
[function_1](["argument"])
\end{lstlisting}

[Describe the output: what was created or returned, and why it is useful.]

\subsection{Realistic Example}
[A longer case that mirrors real use: options, combinations of functions, edge cases. Show the output (table, figure, file tree) as in \cref{fig:output}.]

\begin{figure}[htbp]
    \centering
    \safegraphic[width=0.8\textwidth]{example_duel.pdf}
    \caption{[Output of the realistic example] Example: a figure drawn in R}
    \label{fig:output}
\end{figure}

% =====================================================================
\section{Quality and Validation}\label{sec:quality} % ~10%
[How do you know it works? Unit tests (\texttt{testthat}, number of tests, coverage), \texttt{R CMD check} status, continuous integration, platforms tested (Windows, macOS, Linux). If the package computes results, compare them with a reference implementation or known values.]

% =====================================================================
\section{Limitations and Future Work}\label{sec:limits} % ~10%
[Be honest: what it does not do, when it fails, which platforms are untested. Then the roadmap: 2-3 concrete next features.]

% =====================================================================
\section{Conclusion}\label{sec:conclusion} % ~5%
[Restate the problem and how \texttt{{{package}}} solves it, in two or three sentences. No new information.]

% =====================================================================
\section*{Software Availability}
\begin{itemize}[nosep]
    \item Repository: {{repo_url}}
    \item Version described: [0.1.0], archived at [Zenodo DOI]
    \item License: [MIT]
    \item Requirements: R $\geq$ [4.1], [dependencies]
\end{itemize}

\section*{Acknowledgements}
[People, funding, and the packages this one builds on.]

\newpage
\printbibliography

\newpage
\appendix
\section{Function Reference}\label{app:reference}
[One entry per exported function: signature, arguments, return value. Can be generated from the roxygen documentation.]

\section{Session Information}\label{app:session}
[Output of \texttt{sessionInfo()} for the examples in this paper.]

\end{document}
